\documentclass[UTF8,12pt,a4paper,fontset=fandol]{ctexart}

% 支持从项目根目录或当前笔记目录编译。
\makeatletter
\def\input@path{{../}{../../}}
\makeatother
\usepackage{researchnotes}

\title{克拉默法则的推导}
\author{}
\date{}

\begin{document}
\maketitle

\section{从二元消元看出规律}

考虑方程组 $ax+by=u$、$cx+dy=v$。分别消去 $y$ 和 $x$，得到
\[
  (ad-bc)x=ud-bv,\qquad (ad-bc)y=av-cu.
\]
当 $ad-bc\ne0$ 时，任何解都必须满足
\[
  x=\frac{\begin{vmatrix}u&b\\v&d\end{vmatrix}}
          {\begin{vmatrix}a&b\\c&d\end{vmatrix}},\qquad
  y=\frac{\begin{vmatrix}a&u\\c&v\end{vmatrix}}
          {\begin{vmatrix}a&b\\c&d\end{vmatrix}}.
\]
分母都是系数行列式；求哪个未知数，分子就把它对应的系数列换成常数项列。
\keyterm{克拉默法则}（Cramer's rule）把这一规律推广到任意阶。
下面只用行列式的线性性质、两列（行）相同时行列式为零，以及按列（行）展开公式来推导。

\section{为什么替换一列就能提取一个未知数}

设 $n$ 为正整数，$A=(a_{ij})\in\mathbb F^{n\times n}$，其中
$\mathbb F=\mathbb R$ 或 $\mathbb C$。记 $A$ 的列为 $a_1,\ldots,a_n$，
常数项向量为 $b=(b_1,\ldots,b_n)^{\mathsf T}$，未知向量为
$x=(x_1,\ldots,x_n)^{\mathsf T}$，$\mathsf T$ 表示转置。
方程组 $Ax=b$ 就是
\begin{equation}
  x_1a_1+\cdots+x_na_n=b.
  \label{eq:derivation-column-system}
\end{equation}
记 $D=\det A$，并对 $j=1,\ldots,n$ 定义
\[
  D_j=\det(a_1,\ldots,a_{j-1},b,a_{j+1},\ldots,a_n).
\]
也就是说，$D_j$ 只替换第 $j$ 列，其他列的位置不变。

假设 $x$ 是一个解。将式~\eqref{eq:derivation-column-system} 代入 $D_j$，
再利用行列式对第 $j$ 列的线性性质，得到
\begin{equation}
  D_j=\sum_{k=1}^n x_k
  \det(a_1,\ldots,a_{j-1},a_k,a_{j+1},\ldots,a_n)=x_jD.
  \label{eq:derivation-extract}
\end{equation}
最后一步的原因是：当 $k\ne j$ 时，第 $j$ 列与原来的第 $k$ 列相同，
该项行列式为零；当 $k=j$ 时，矩阵恰好恢复为 $A$。
因此，\keyterm{替换列再作线性展开，会消去其余未知数，只留下所求未知数乘以 $D$}。

\begin{theorem}[克拉默法则]
\label{thm:derivation-cramer}
若 $D\ne0$，则对任意 $b\in\mathbb F^n$，方程组 $Ax=b$ 有唯一解，且
\begin{equation}
  x_j=\frac{D_j}{D},\qquad j=1,\ldots,n.
  \label{eq:derivation-cramer}
\end{equation}
\end{theorem}

式~\eqref{eq:derivation-extract} 已说明任何解都只能取上述值，因而解至多有一个。
但推导时假设了解存在，还须验证这些比值确实满足原方程组。

\section{补全证明：这些比值确实构成解}

\begin{proof}[定理~\ref{thm:derivation-cramer} 的存在性证明]
记 $C_{ij}$ 为 $a_{ij}$ 的\keyterm{代数余子式}（cofactor），即删去 $A$ 的第 $i$ 行、
第 $j$ 列后所得行列式乘以 $(-1)^{i+j}$；$n=1$ 时约定 $C_{11}=1$。
沿替换后的第 $j$ 列展开，因为删去这一列后剩余元素没有改变，故
\[
  D_j=\sum_{i=1}^n b_iC_{ij}.
\]
对任意 $r,i\in\{1,\ldots,n\}$，按行展开又给出
\begin{equation}
  \sum_{j=1}^n a_{rj}C_{ij}
  =\begin{cases}D,&r=i,\\0,&r\ne i.\end{cases}
  \label{eq:derivation-cofactor}
\end{equation}
当 $r=i$ 时，这是 $D$ 按第 $i$ 行展开。
当 $r\ne i$ 时，把 $A$ 的第 $i$ 行换成第 $r$ 行，再按第 $i$ 行展开，
就得到左端；删去该行后子式不变，而替换后的矩阵有两行相同，所以结果为零。

令 $\widehat x_j=D_j/D$。对原方程组的每一行 $r$，有
\[
  \sum_{j=1}^n a_{rj}\widehat x_j
  =\frac1D\sum_{j=1}^n a_{rj}\sum_{i=1}^n b_iC_{ij}
  =\frac1D\sum_{i=1}^n b_i\sum_{j=1}^n a_{rj}C_{ij}
  =b_r.
\]
这里仅交换了有限求和，再使用式~\eqref{eq:derivation-cofactor}。
因此 $A\widehat x=b$，解确实存在；结合前面得到的唯一性，证明完成。
\end{proof}

条件 $D\ne0$ 用于把 $D_j=x_jD$ 除以 $D$。
当 $D=0$ 时，任何解仍须满足 $D_j=x_jD=0$，所以只要某个 $D_j\ne0$ 就无解；
若所有 $D_j=0$，则不能用这些比值求解，应另用消元法判断。

\end{document}
